%MIT OpenCourseWare: https://ocw.mit.edu
%RES.18-012 Algebra II Student Notes, Spring 2022
%License: Creative Commons BY-NC-SA 
%For information about citing these materials or our Terms of Use, visit: https://ocw.mit.edu/terms.

% \rhead{\rightmark}
\lhead{Lecture 20: Modules and Presentation Matrices}

\section{Modules and Presentation Matrices}

\subsection{Review --- Definition of Modules}

Last class, we defined modules over commutative rings --- we've also seen a few examples over noncommutative rings, but from now we'll stick to commutative ones. 

\begin{definition}
A \textbf{module} $M$ over a ring $R$ is an abelian group together with an action of $R$ --- each $r \in R$ acts by a map $r : m \mapsto r(m)$ (we often denote $r(m)$ by $rm$), satisfying certain axioms. 
\end{definition}

In some sense, modules can be similar to vector spaces: 

\begin{example}
The free module of rank $n$ is $M = R^n$, consisting of $n$-tuples of elements in $R$ (where addition and the $R$-action are performed componentwise). 
\end{example}

But there are many other examples of modules over most rings, and we'll now discuss how to describe more general modules. 

\subsection{Generators and Relations}

If $M$ is a module, the elements $a_1, \ldots, a_n \in M$ form a \textbf{system of generators} if every $x \in M$ has the form \[x = \sum r_i a_i\] for some $r_i \in R.$ 

Choosing any $n$ elements $a_1, \ldots, a_n \in M$ provides a homomorphism of modules 
$\varphi : R^n \to M$ (where we send $(r_1, \ldots, r_n) \mapsto r_1a_1 + \cdots + r_na_n$). Then $a_1$, \ldots, $a_n$ are generators if and only if $\varphi$ is onto. 

In the analogy with vector spaces, a system of generators is a set of vectors that \emph{span} the vector space. But in vector spaces, if we have a set of vectors which span the space, we can always find a subset which forms a basis --- we can drop some of the $a_i$ to make the presentation $x = \sum r_ia_i$ be \emph{unique}. This is not true in general. 

Stating that the presentation $x = \sum r_ia_i$ is unique is equivalent to stating that $\ker \varphi = 0$ (given two presentations, we could subtract them to get an element of the kernel). Then $\varphi : R^n \to M$ is actually an isomorphism, which means $M$ is free. But this is often \emph{not} the case --- there are usually many $R$-modules which are not free. 

\begin{question}
Does the system of generators have to be finite?
\end{question}

\begin{ans}
Not necessarily; as a silly example, we could take \emph{all} elements of $M$ as our system of generators. In later classes, we'll discuss ways to sometimes show that we \emph{can} always find a finite system of generators; but for today, we'll just \emph{assume} that we can. 
\end{ans}

\begin{definition}
If a finite set of generators exists, then $M$ is called \textbf{finitely generated.} 
\end{definition}

\subsection{Presentation Matrices}

Most of the time, we won't be lucky enough that $\ker \varphi = 0$. But in general, by the homomorphism theorem, we can write \[M \cong R^n/\ker \varphi.\] Now $\ker \varphi$ is \emph{also} a module, so we can describe it further. 

\begin{definition}
If $\ker \varphi$ is also finitely generated (for some surjective homomorphism $\varphi : R^n \to M$), then we say $M$ is \textbf{finitely presented}. 
\end{definition}

We'll see later that for a large class of rings, \emph{every} module which is finitely generated is also finitely presented; but for now, we'll assume that $M$ is finitely presented as well. 

Then we can choose a set of generators $b_1$, \ldots, $b_m$ for $\ker \varphi$. Each $b_i$ can be thought of as a column vector (since it's an element of $R^n$). So we can write down the $n \times m$ marix \[B = \begin{bmatrix} \vert & \vert &  & \vert \\ b_1 & b_2 & \cdots & b_m \\ \vert & \vert & & \vert \end{bmatrix},\] called the \textbf{presentation matrix}. Now since we've fixed generators for $\ker \varphi$, we have an onto map $\psi : R^m \to \ker \varphi$. Equivalently, we can think of $\psi$ as a map $\psi : R^m \to R^n$ whose image is exactly $\ker \varphi$. As in the case of vector spaces in linear algebra, this map can explicitly described as $\psi : x \mapsto Bx$ (where $x \in R^m$), and $\im \psi = BR^m$. This means we can write \[M \cong R^n/BR^m\] (where $BR^m$ is the span of the column vectors of the presentation matrix $B$). 

% By the homomorphism theorem, \todo{what theorem} $M \cong R^n / \ker \varphi.$ If $\ker \varphi$ is also finitely generated, then $M$ is finitely presented. In this case, it is possible to choose a set of generators $b_1, \cdots, b_m \in \ker \varphi.$ Each $b_i \in R^n$ can be thought of as a column vector $\textbf{b}_i$. Then, let the presentation matrix $B$ be the $n \times m$ matrix \[B = [\textbf{b}_1 \textbf{b}_2 \cdots \textbf{b}_m].\] 

% We have 
% \begin{align*}
%     \psi: &R^m \rto \ker \varphi \\
%     \psi: &R^m \rto R^n,
% \end{align*}

% where $\im(\psi) = \ker(\varphi).$ The mapping is such that $\psi(x) = B \cdot \textbf{x}$, where $x \in R^m$ and $\textbf{x}$ is a column vector. Then $M = R^n/B\cdot R^m,$ where $B \cdot R^m$ is the span of the column vectors of the presentation matrix $B.$

\subsection{Classifying Modules}

The rest of the lecture focuses on classifying finitely generated modules over a Euclidean domain. Since abelian groups are $\ZZ$-modules, this will also give us a classification of finitely generated groups. 

For a given module $M$, the presentation is not at all unique. 

\begin{example}
Let $R = \ZZ$ and $M = \ZZ/5\ZZ$. An obvious presentation of $M$ is to choose one generator $1$, and the generator $5$ for the kernel. In this case, $B = \begin{bmatrix} 5 \end{bmatrix}$. 
But there are other presentations as well. For example, we can choose two generators for $M$, and the presentation matrix \[B = \begin{bmatrix} 2 & 1 \\ 1 & 3 \end{bmatrix}.\] To see that $\ZZ^2/B\ZZ^2$ is again $\ZZ/5\ZZ$, note that the sublattice $L \subset \ZZ^2$ spanned by $(2, 1)^t$ and $(1, 3)^t$ has index $\lvert \det(B) \rvert = 5$, which means the quotient $\ZZ^2/L$ has five elements. But the only group of five elements is $\ZZ/5\ZZ$, so this construction indeed gives another presentation of $\ZZ/5\ZZ$. 

% An obvious presentation is $B = [5].$\todo{explain this one again} A different presentation is $B = \begin{pmatrix}
% 2 & 1 \\ 1 & 3
% \end{pmatrix}$, for which $\det(B) = 5.$ Again, $M = \ZZ/5\ZZ.$ The sublattice $L$ in the standard lattice $\ZZ^2$ spanned by the column vectors has index $|\det(B)| = 5,$ so the quotient $\ZZ^2 / L $ has five elements, and the only abelian group of 5 elements is $\ZZ/5.$
\end{example}

Our goal is to more systematically understand how to understand the module given a presentation. Here, the analogy between vector spaces and modules will be useful --- similarly to in linear algebra, we'll start with a matrix and try to perform elementary operations on the matrix which don't change the module. We'll then use these operations to write the matrix in a simpler form, from where it's easy to understand the module. 

% Since vector spaces are an example of modules, our understanding of linear algebra will again come in useful. \todo{write down what he said here}

\subsubsection{Elementary Row and Column Operations}

We'll use the notation $M_B$ to refer to the module $M$ produced by a presentation matrix $B$. 

In linear algebra, we saw the elementary column operations for matrices over a field. We can define elementary column operations on matrices over a \emph{ring} in a similar way:
\begin{enumerate}
    \item Multiplying a column by a unit (an invertible element) --- note that in the case of a field, we could multiply by \emph{any} nonzero element (because any nonzero element has an inverse), but it was important that we were able to invert the multiplication; so here we restrict the definition to units. 
    \item Adding an arbitrary multiple of one column to another column. 
    \item For convenience, we can also include swapping two columns; but in fact, this can be obtained as a combination of the first two operations. 
\end{enumerate}

These operations \emph{do not} change the span of the columns. (This can be verified the same way as for matrices over a field.) So if $B'$ is obtained from $B$ by elementary column operations, then we have $BR^m = B'R^m$. 

Another useful way to think of this is in terms of matrix multiplication. In other words, if $B'$ is obtained from $B$ by elementary column operations, then we have $B' = B\cdot C$, where $C$ is an $m \times m$ \emph{invertible} matrix (it's easy to see that all the elementary column operations are invertible). This means $C \in \mathrm{GL}_m(R)$ (the group of invertible matrices wth entries in $R$). Note that over a ring, for a matrix $C$ to be invertible, $\det(C)$ must be a \emph{unit} --- it's not enough to require the determinant to be nonzero. (In fact, the converse is also true, but requires more work.)

Then we have $B'R^m = BCR^m$. But since $C$ is invertible, the map defined by $C$ is an isomorphism, so $CR^m = R^m$. So this is an alternate way of making it clear that $B'R^m = BR^m$. 

% So $BR^m = B' R^m$ if $B'$ is obtained from $B$ by these column operations. This is true because if $B'$ is obtained from $B$ by column operations, then $B' = B\cdot C$ where $C$ is an $m \times m$ invertible matrix. Note that $C \in GL_m(R)$, the group of invertible matrices with entries in $R,$ and for $C$ to be invertible, $\det(C)$ must be a \emph{unit} in $R;$ it is not enough to require that the determinant of $C$ be nonzero.\footnote{The converse is also true, but requires proof.} Then $B' R^m = B(CR^m)$, and since $C$ is invertible, $CR^m = R^m,$ and so $B'R^M = BR^m,$ and so they have the same span.

The elementary \emph{row} operations are defined analogously. 

When we apply elementary row operations to a matrix $B$, we produce a matrix $B'$ which is a presentation matrix for an \emph{isomorphic} module --- in order to explain why, we'll again think in terms of matrix multiplication. We can write $B' = CB$, where $C \in \mathrm{GL}_n(R)$. We then want to produce an isomorphism between $R^n/BR^m$ and $R^n/CBR^m$. But that isomorphism is just given by multiplication by $C$. More explicitly, we can draw a commutative diagram:


% , and when applied to $B$ produce a matrix $B'$ which is a presentation for an \emph{isomorphic} module: $M_B \cong M_{B'}.$ In this case, $B' = CB,$ where $C \in GL_n(R).$ We want an isomorphism between $R^n/B\cdot R^m$ and $R^n/CB \cdot R^m,$ which is given by $C.$ The following diagram commutes, where the top arrow is an isomorphism:

\begin{center}
    % https://tikzcd.yichuanshen.de/#N4Igdg9gJgpgziAXAbVABwnAlgFyxMJZABgBpiBdUkANwEMAbAVxiRACUA9MAegCEuAWxABfUuky58hFAEZyVWoxZsuvAMJ8AOloDGUCDgAEQ0eJAZseAkTKzF9Zq0QduZiVelF596o5UuaqKKMFAA5vBEoABmAE4QwogATNQ4EEgAzNQMdABGMAwACpLWMiCxWGEAFjggfsrOIOruIHEJSGQgaUiyYjHxiSld6YjEfa0DmakjvRQiQA
\begin{tikzcd}
R^n/BR^m \arrow[r]            & R^n/CB\cdot R^m \\
R^n \arrow[r, "C"'] \arrow[u] & R^n \arrow[u]  
\end{tikzcd}
\end{center}

The map $x \mapsto Cx$ is an isomorphism $R^n \to R^n$. But by definition, $x \in BR^n$ if and only if $Cx \in CBR^m$. So if we restrict our isomorphism to $BR^m$, then it restricts to an isomorphism between $BR^m$ and $CBR^m$; and therefore to an isomorphism between the quotients $R^n/BR^m$ and $R^n/CB\cdot R^m$. %, where $\overline{x} \in BR^m \equiv C\overline{x} \in CB \cdot R^m$, so it restricts to an isomorphism $BR^m \rto B'R^m$ and hence an isomorphism of quotients. 

So the row and column operations don't change our module (up to isomorphism). 

\subsubsection{Smith Normal Form}

Using the row and column operations, it's possible to write any presentation matrix in a much simpler form. (We'll primarily focus on the case $R = \ZZ$, but this holds for any Euclidean domain.)

\begin{theorem} \label{Smith normal form}
Every $n \times m$ matrix over a Euclidean domain $R$ can be reduced by elementary row and column operations to a matrix in \emph{Smith normal form} --- if we let $B = (b_{ij})$, then we have $b_{ij} = 0$ for all $i \neq j$, and $b_{11} \mid b_{22} \mid b_{33} \mid \cdots$.  %where $b_{ij} = 0$ if $i \neq j,$ and letting $d_i = b_{ii}$, $d_i$ divides $d_{k}$ if $i < k.$ \todo{explain the smith normal form more}
\end{theorem}

So a matrix in Smith normal form looks like \[\begin{bmatrix} d_1 \\ & d_2 \\ & & d_3 \\ & & & \ddots \\ & & & & & d_k & & & \end{bmatrix}\] (where all entries not shown are $0$'s). 

We'll discuss the proof next class. It combines two ideas --- the method of using Gaussian elimination to solve systems of equations over a \emph{field} (which involves reducing matrices to a simpler form using row and column operations as well), and the Euclidean algorithm. (We've previously discussed Euclidean domains in the context of them being PIDs, but it turns out that here, having an effective way of computing the gcd using division with remainder will be very useful. The theorem is also true for PIDs, with a bit of modification (and a different proof); but all the examples we'll be interested in are Euclidean domains. 

% A matrix in Smith normal form is diagonal. The proof combines similar reduction by Gaussian elimination as for $R = F$ and the Euclidean algorithm. In fact, this is also true when $R$ is a PID. 

\begin{corollary}
Every finitely presented module over a Euclidean domain is isomorphic to a direct sum of \emph{cyclic} modules (modules which are generated by one element) --- we can write \[M \cong R^a \oplus R/(d_1) \oplus R/(d_2) \oplus \cdots \oplus R/(d_k),\] where we additionally have $d_1 \mid d_2 \mid \cdots \mid d_k$. 
\end{corollary}

Later we'll see that any finitely generated module over a Euclidean domain is also finitely presented; this will mean that our statement actually holds for all finitely \emph{generated} modules. 

This corollary is clear from the theorem:

\begin{proof}
Using Theorem \ref{Smith normal form}, we can rewrite the presentation matrix $B$ to be diagonal, with diagonal $d_1 \mid d_2 \mid \cdots \mid d_k$, where $k = \min(m, n)$. In this case, when we take a column vector in $R^m$ and multiply by $B$, we simply scale each coordinate by the corresponding $d_i$ --- so when we quotient out by $BR^m$, this coordinate becomes $R/(d_i)$, and we get \[M_B = \bigoplus R/(d_i) \oplus R^a.\] (The extra free factor comes from rows with no entries --- either because $d_i = 0$ or because $n > m$ --- since such rows correspond to coordinates in $R^n$ where we're not quotienting out by anything.)
\end{proof}

% The corollary is clear from the theorem because the image of a diagonal matrix scales each coordinate separately by a number. For a vector to lie in the image, every coordinate is scaled by the corresponding number. 

% \begin{proof}

% Since the $n \times m$ matrix $B$ is diagonal with $d_1, \cdots, d_n$ on the diagonal, $M_B = \oplus_{i = 1}^n (R/d_i).$
% \end{proof}

In fact, this classification can be used to understand the subgroups of lattices as well (which came up when we studied factorization in quadratic number fields). This theorem implies that to describe a subgroup, we can always choose a basis for the lattice, and simply scale these basis vectors by numbers to get a basis for the subgroup. 

\begin{question}
What does it mean for a module to be cyclic --- is the free module cyclic? 
\end{question}

\begin{ans}
A module is cyclic if it's generated by one element. The free module is cyclic --- it's generated by one element with \emph{no} relations. Meanwhile, $R/d$ is generated by one element $x$, with the relation that $dx = 0$. 
\end{ans}

% \todo{he has some explanation here}
